% sections/01-connection-and-curvature.tex
\section{The connection, the second fundamental form, and the Gauss equation}
\label{sec: connection-curvature}

This section and the next supply the apparatus, the curvature of the
configuration group and its size, which Sections~\ref{sec: dichotomy}
and~\ref{sec: blowup} turn on the singular set.

Let $\Lsig$ be the space of smooth divergence-free vector fields on $\T^3$ with
the $L^2$ inner product $g(v,w) = \int_{\T^3} v\cdot w\,\dd x$. The Leray
projection $\Leray$ is the orthogonal projection of $L^2$ onto $\Lsig$, and
$\Helm = I - \Leray = \nabla\Delta^{-1}\diverg$ is the complementary projection
onto gradients. The group $\Diffop(\T^3)$ of all diffeomorphisms carries a flat
$L^2$ metric, whose geodesics are the pressureless straight-line motions; its
volume-preserving subgroup $\SDiff(\T^3)$ is a submanifold whose induced
Levi-Civita connection on $\Lsig$ is the projection of the flat
connection~\cite{ebin-1970-gro-dif-mot, misiolek-1993-sta-flo-ide,
arnold-1998-top-met-hyd},
\begin{equation}\label{eq: connection-2}
  \conn_u v = \Leray\bigl[(u\cdot\nabla)v\bigr], \qquad u, v \in \Lsig.
\end{equation}
This is the metric connection in the convention of Misio\l{}ek and of Ebin and
Marsden; the Gauss equation of Proposition~\ref{prop: gauss} confirms that the
sectional curvature it produces is Arnold's.

\subsection{The second fundamental form}

The normal part of the flat derivative is the second fundamental form of the
constraint.

\begin{lemma}\label{lem: second-fundamental}
The second fundamental form of $\SDiff(\T^3)$ in $\Diffop(\T^3)$ is
\begin{equation}\label{eq: II}
  \sff(a,b) = \Helm\bigl[(a\cdot\nabla)b\bigr] = \nabla\phi_{ab},
  \qquad \Delta\phi_{ab} = \partial_i a_j\,\partial_j b_i,
\end{equation}
for $a, b \in \Lsig$. It is symmetric, and on the diagonal it is the pressure
gradient, $\sff(u,u) = -\nabla p$.
\end{lemma}

\begin{proof}
The flat derivative of $b$ along $a$ is $(a\cdot\nabla)b$; its normal
(gradient) part is $\Helm[(a\cdot\nabla)b] = \nabla\Delta^{-1}\diverg[(a\cdot\nabla)b]$,
and $\diverg[(a\cdot\nabla)b] = \partial_i a_j\,\partial_j b_i$ since $a$ is
divergence-free. Symmetry follows because the Lie bracket
$[a,b] = (a\cdot\nabla)b - (b\cdot\nabla)a$ is divergence-free, so
$\Helm[(a\cdot\nabla)b] = \Helm[(b\cdot\nabla)a]$. For $a = b = u$ the source is
$\partial_i u_j\,\partial_j u_i = q = -\Delta p$ by~\eqref{eq: pressure-poisson},
hence $\phi_{uu} = -p$ and $\sff(u,u) = -\nabla p$.
\end{proof}

\subsection{Curvature and the Gauss equation}

The curvature of~\eqref{eq: connection-2},
\[
  \curv(u,v)w = \conn_u\conn_v w - \conn_v\conn_u w - \conn_{[u,v]} w,
\]
expands, for $u, v, w \in \Lsig$, as
\begin{equation}\label{eq: curvature-formula}
  \curv(u,v)w
  = \Leray\!\bigl[(u\cdot\nabla)\Leray[(v\cdot\nabla)w]\bigr]
  - \Leray\!\bigl[(v\cdot\nabla)\Leray[(u\cdot\nabla)w]\bigr]
  - \Leray\!\bigl[([u,v]\cdot\nabla)w\bigr].
\end{equation}
Since the ambient group is flat, the curvature is governed entirely by the second
fundamental form through the Gauss equation. This is the content of the main
identity.

\begin{theorem}\label{thm: main}
Let $u$ be smooth and divergence-free
on $\T^3$. The flat ambient part of the curvature vanishes identically,
\begin{equation}\label{eq: local-zero}
  (u\cdot\nabla)(v\cdot\nabla)u - (v\cdot\nabla)(u\cdot\nabla)u
  - ([u,v]\cdot\nabla)u = 0,
\end{equation}
and for every divergence-free $v$ the curvature operator is the Leray-projected
pressure Hessian,
\begin{equation}\label{eq: reduced}
  \curv(u,v)u = -\Leray\bigl[(\nabla^2 p)\,v\bigr]
              - \Leray\bigl[(\nabla^2 q_v)\,u\bigr],
\end{equation}
with $q_v$ the mixed pressure $\Delta q_v = \partial_i v_j\,\partial_j u_i$.
\end{theorem}

Equation~\eqref{eq: reduced} names the curvature as the pressure Hessian
$\nabla^2 p = -\Riesz q$ acting on $v$, together with a mixed-pressure term
quadratic in the velocity gradient. The sectional form of the same identity is
Arnold's, recovered here from the second fundamental form.

\begin{proposition}\label{prop: gauss}
For $u, v \in \Lsig$ the sectional
curvature satisfies the Gauss equation
\begin{equation}\label{eq: gauss}
  g\bigl(\curv(u,v)v, u\bigr)
  = \bigl\langle \sff(u,u), \sff(v,v) \bigr\rangle_{L^2}
  - \norm{\sff(u,v)}_{L^2}^2,
\end{equation}
which is Arnold's sectional curvature formula with the pressure term
$-\norm{\sff(u,v)}_{L^2}^2$ and remainder
$B(u,v) = \langle \sff(u,u), \sff(v,v)\rangle_{L^2}$.
\end{proposition}

\begin{proof}
The ambient group $\Diffop(\T^3)$ is flat, so the Gauss equation for the
submanifold $\SDiff(\T^3)$ reads
$g(\curv(u,v)v,u) = \langle\sff(u,u),\sff(v,v)\rangle - \langle\sff(u,v),\sff(v,u)\rangle$,
and $\sff$ is symmetric by Lemma~\ref{lem: second-fundamental}. With
$\sff(u,u) = -\nabla p$ this is Arnold's formula; the pressure energy
$-\norm{\sff(u,v)}^2$ is its leading term and $B = \langle\sff(u,u),\sff(v,v)\rangle$
the remainder that vanishes on the Fourier-mode pairs of~\cite{arnold-1998-top-met-hyd}.
\end{proof}
